\begin{abstract}

通过阳极氧化的方法在金属表面形成纳米结构的氧化层已经被广泛用于光催化和储能材料等领域，但是阳极氧化物生长机理的研究仍缺乏系统性和连续性。本课题基于目前研究较完善的阳极氧化TiO2纳米管的生长机理，通过定量分析离子电流和电子电流等与多变的阳极氧化纳米形貌的关系，研究和建立多孔阳极氧化Sn的生长模型，并指出了两者的不同点，同时……

\keywords{阳极氧化TiO2纳米管，多孔阳极氧化锡}

\end{abstract}




\begin{englishabstract}

Nanostructured oxide layer formed on metal surface by anodic oxidation method has been widely used in photocatalysis and energy storage materials, but the growth mechanism of anodic oxide is still not systematic and continuous. This research is based on the growth mechanism of anodized TiO2 nanotubes which is well studied at present, and through quantitative analysis of ionic and electronic current and the relationship with variable anodized nano morphologies, the growth mechanism of porous anodic Sn oxide was studied and compared with that of anodized TiO2 nanotubes. Meanwhile…… 

\englishkeywords{anodized TiO2 nanotubes, porous anodic Sn oxide}

\end{englishabstract}